%
% 1-elementar.tex -- Elementare Funktionen
%
% (c) 2026 Prof Dr Andreas Müller
%
\section{Elementare Funktionen
\label{buch:intalg:section:elementar}}
\kopfrechts{Elementare Funktionen}%
Der Benutzer eines Computeralgebrasystems erwartet, dass er
als Integranden Funktionen wie
\[
\!\sqrt{ax^2+bx+c},\quad
\sin(x-\pi),\quad
\sqrt{\tan x},\quad
e^{-x^2}
\]
eingeben kann und als Output eine Stammfunktion in einer ähnlichen
Form erhalten kann.
Zum Beispiel liefert Maxima für die dritte Funktion die Stammfunktion
\begin{align*}
\int\!\sqrt{\tan x}\,dx
&=
-\frac1{2^{\frac32}}
\Bigl(
\log \bigl(\tan x+\!\sqrt{2}\sqrt{\tan x}+1\bigr)
-
\log \bigl(\tan x-\!\sqrt{2}\sqrt{\tan x}+1\bigr)
\Bigr)
\\
&\qquad
+\frac1{\!\sqrt{2}}
\Bigl(
\arctan \bigl( \!\sqrt{2}\sqrt{\tan x}+1 \bigr)
+
\arctan \bigl( \!\sqrt{2}\sqrt{\tan x}-1 \bigr)
\Bigr).
\end{align*}
Natürlich ist es recht anspruchsvoll, auf diese Lösung zu kommen.
Allein durch Anwendung heuristischer Integrationsregeln kommt
man wohl kaum darauf.
Es braucht also einen algebraischen Prozess, der diese sehr 
komplizierte Stammfunktion schrittweise aufbauen kann.

In der Stammfunktion kommt nicht überraschend die
Funktion $\!\sqrt{\tan x}$
wieder vor, aber auch Polynome in dieser Funktion sowie Logarithmus
und Arkustangens-Funktionen, mit solchen Polynomen als Argument.
\index{arctan}%
Schreibt man $\vartheta(x)=\!\sqrt{\tan x}$, dann wird die Stammfunktion
zu
\begin{align*}
\int\vartheta(x)\,dx
&=
-\frac1{2\!\sqrt{2}}
\Bigl(
\log \bigl(\vartheta(x)^2 +\!\sqrt{2} \vartheta(x) + 1\bigr)
-
\log \bigl(\vartheta(x)^2 -\!\sqrt{2} \vartheta(x) + 1\bigr)
\Bigr)
\\
&\qquad
+\frac1{\!\sqrt{2}}
\Bigl(
\arctan \bigl( \!\sqrt{2}\vartheta(x) + 1 \bigr)
+
\arctan \bigl( \!\sqrt{2}\vartheta(x) - 1 \bigr)
\Bigr).
\end{align*}
Um diese Funktion aufzubauen, wird zusätzlich zur Funktion $\vartheta(x)$
zunächst die Konstante $\!\sqrt{2}$ benötigt.
Damit lassen sich dann die Polynome
\[
\vartheta(x)^2 + \!\sqrt{2}\vartheta(x) + 1,
\quad
\vartheta(x)^2 - \!\sqrt{2}\vartheta(x) + 1,
\quad
\!\sqrt{2}\vartheta(x)+1,
\quad
\!\sqrt{2}\vartheta(x)-1
\]
in $\vartheta(x)$ bilden.
Im nächsten Schritt muss dann der Logarithmus bzw.~die Arkustangensfunktion
angewendet werden.
Schliesslich kann die Stammfunktion als Linearkombination gebildet werden.

Jahrelange Beschäftigung mit algebraischen Umformungen hat die Intuition
gebildet, die in der obenstehenden Analyse der Struktur der Formel
angewendet worden ist.
Um auf diesen Ideen aufbauend einen Integrationsalgorithmus herstellen
zu können, müssen wir die einzelnen Schritt durch wohldefinierte
algebraische Konstruktionen ersetzen.

%
% Funktionenkörper
%
\subsection{Funktionenkörper}
In Kapitel~\ref{chapter:ratint} wurden das Problem studiert,
rationale Funktionen zu integrieren.
Die meisten Operationen bei der Lösung dieses Problems ergaben
wieder rationale Funktionen.
Nur ganz am Schluss traten Terme auf, deren Integrale mit dem
Logarithmus eines Polynoms berechnet werden können.
Wir brauchen also eine Menge von Funktionen, in der alle arithmetischen
Operationen durchgeführt werden können, also einen Körper von
Funktionen.

%
% Koeffizientenkörper
%
\subsubsection{Koeffizientenkörper}
In Abschnitt~\ref{buch:ratint:section:bernoulli} haben wir als
Funktionenkörper den Körper $\mathbb{C}(x)$ verwendet, dann
aber später eingesehen, dass der Körper $\mathbb{C}$ zu gross ist.
Wenn man mit irrationalen Zahlen algebraisch rechnen muss, muss
man deren algebraischen Eigenschaften kennen.
Die Nullstellen eines Polynoms, die für die Partialbruchzerlegung
benötigt wurden, sind algebraische Element über einem Körper, der
alle Koeffizienten des Nennerpolynoms enthält.
Der Ausgangskörper für die Konstruktion der rationalen Funktionen
entsteht also durch Hinzufügen weiterer Konstanten zu den rationalen
Zahlen.
Die zusätzlichen Konstanten können algebraisch sein, wie die
Wurzel $\!\sqrt{2}$, oder transzendent, wie $\pi$, für die es keine
polynomielle Relation gibt.
Als Zahlenkörper, aus dem die Koeffizienten der Polynome stammen
muss daher immer in Erweiterung von $\mathbb{Q}$ verwendet werden,
die durch Hinzufügen endlich vieler Elemente $c_1,\dots,c_n$ entstehen.
In der Kette
\[
\underbrace{\mathbb{Q}}_{\displaystyle = K_0}
\subset
\underbrace{\mathbb{Q}(c_1)}_{\displaystyle = K_1}
\subset
\underbrace{\mathbb{Q}(c_1,c_2)}_{\displaystyle = K_2}
\subset
\underbrace{\mathbb{Q}(c_1,c_2,c_3)}_{\displaystyle = K_3}
\subset
\ldots
%\subset
%\underbrace{\mathbb{Q}(c_1,c_2,c_3,\ldots,c_{n-1})}_{\displaystyle = K_{n-1}}
\subset
\underbrace{\mathbb{Q}(c_1,c_2,c_3,\ldots,c_{n-1},c_n)}_{\displaystyle = K_n}
=
K.
\]
Jedes Element $c_k$ ist transzendent oder algebraisch ist über
dem vorangegangenen Körper.
Falls $c_k$ algebraisch ist, dann gibt es ein Minimalpolynom
$m_k(x) \in K_{k-1}[x]$, welches $m_k(c_k)=0$.

Zu den transzendenten Elementen in $K$ gehören insbesondere auch
Parameter in den Ausdrücken.
Wenn ein CAS den quadratischen Ausdruck $ax^2+bx+c$ integrieren soll,
ohne den Koeffizienten $a$, $b$ und $c$ feste Werte zuzuordnen, dann
müssen diese Elemente als transzendet über jedem zugrundeliegenden
Körper betrachtet werden.
Damit sagt man, dass keine weiteren Beziehungen zwischen diesen
Koeffizienten bekannt sind.

%
% Das Funktionsargument
%
\subsubsection{Das Funktionsargument}
In einem Funktionsausdruck spielt ein Symbol die spezielle Rolle
des Funktionsarguments.
Aus algebraischer Sicht verhält sich der Platzhalter für das
Funktionsargument nicht anders wie jeder andere Parameter.
Im Ausdruck $f(x) = ax^2 + bx + c$ werden die Zeichen $a$, $b$ und $c$
wie gewöhnliche Zahlen behandelt, während das Symbol $x$ für das
Funktionsargument steht, für das man nach der Intuition, die man
beim ersten Kontakt mit dem Funktionsbegriff kennenlernt, beliebige
Werte einsetzen kann.
Aus algebraischer Sicht ist $x$ einfach nur ein weiteres transzendentes
Element, um das der Koeffizientenkörper erweitert wird.
Eine besondere Bedeutung erhält $x$ erst, wenn wir, wie das im
nächsten Abschnitt geschehen soll, mit der Ableitung die Konstanten
von der Variablen $x$ unterscheiden können.

\begin{beispiel}
Die rationale Funktion
\[
f(x) 
=
\frac{ax^2 + \!\sqrt{2}x - \sqrt{\pi}}{x^2+\pi}
\]
hat im Zählerpolynom die drei Koeffizienten $a$, $\!\sqrt{2}$ und
$\!\sqrt{\pi}$, im Nenner kommt ausserdem der Koeffizient $\pi$ hinzu.
Die Elemente $a$ und $\pi$ sind transzendent über den rationalen
Zahlen.
Die Elemente
$\!\sqrt{2}$
und
$\!\sqrt{\pi}$
entstehen
durch algebraische Erweiterung von $\mathbb{Q}$ mit dem
Minimalpolynom $m(x)=x^2-2$ bzw.~von $\mathbb{Q}(\pi)$ mit
dem Minimalpolynom $m(x)=x^2-\pi$.
Die Funktion ist daher ein Element im Körper der rationalen
Funktionen
\[
f(x)
\in
\mathbb{Q}(\!\sqrt{2},\pi,\!\sqrt{\pi},a,x).
\]
Man könnte auch versuchen, $\!\sqrt{\pi}$ als transzendentes Element
über $\mathbb{Q}$ zu betrachten.
Dann müsste man aber das Element $\pi$ als das Quadrat von $\!\sqrt{\pi}$
betrachten.
In diesem Fall wäre es zweckmässiger, die Funktion als
\[
f(x)
=
\frac{ax^2+\!\sqrt{2}-c}{x^2+c^2}
\]
mit einem transzendenten Element $c$ zu schreiben.
Diese Version macht es aber schwieriger, eine Beziehung zwischen $\pi$ 
und den trigonometrischen Funktionen herzustellen.
\end{beispiel}

%
% Differentialkörper von Funktionen
%
\subsection{Differentialkörper von Funktionen}
In einem Funktionsausdruck für eine differenzierbare Funktion spielt
unter das Funktionsargument unter allen Symbolen 
eine besondere Rolle.
Wir gehen davon aus, dass man nach dem Funktionsargument ableiten
kann, weil es ``variabel'' ist.
Nach den Parametern gedenkt man dagegen nicht abzuleiten, da man
diese als ``konstant'' ansehen möchte.

%
% Derivation
%
\subsubsection{Derivation}
Im rein algebraischen Framework eines Funktionenkörpers ist es nicht
sinnvoll, die Definition der Ableitung als Grenzwert
\[
f'(z)
=
\lim_{h\to 0}
\frac{f(z+h)-f(z)}{h}
\]
eines Differenzenquotienten zu verwenden.
Um den Grenzwert zu definieren, wird ein Ordnungsrelation gebraucht,
ein Konzept, welches sich nicht mit rein algebraischen Eigenschaften
erfassen lässt.
Für den Grenzwert wird ein Zahlenkörper gebraucht, in dem alle
Grenzwerte existieren, was in $\mathbb{Q}$ nicht gegeben ist.
Der einzige Körper, der vollständig wäre, ist $\mathbb{R}$, der aber
im Computer nicht exakt dargestellt werden kann und daher nicht als
Basis eines CAS definiert werden kann.

Als alternativer Ansatz für die Definition der können die einfachsten
bekannten Rechenregeln der Ableitung verwendet werden.
Die Ableitung ist zum Beispiel eine lineare Abbildung.
Dies allein reicht jedoch nicht, weil die Ableitung auch für 
Potenzen $z^n$ oder für Produkte von Funktionen angewendet werden
soll.
Die Linearität kann darüber nichts aussagen.
Daher braucht es eine zusätzlich Regel für Produkte, die natürlich
die Produktregel für dei Ableitung wiedergeben soll.

\begin{definition}[Derivation]
Eine \emph{Derivation} in einem Funktionenkörper $K$ ist eine 
lineare Abbildung $D\colon K\to K$, die zusätzlich die Produktregel
\[
D(fg) = (Df)g + f(Dg)\qquad\forall f,g\in K
\]
erfüllt.
\index{Derivation}%
\end{definition}

In der Definition steht nichts davon, nach welcher Variablen da
abgeleitet werden soll.
In der Analysis verwendet man die Notation
\[
\frac{\partial}{\partial a},\quad
\frac{\partial}{\partial b},\quad
\frac{\partial}{\partial c},\quad
\frac{\partial}{\partial x},\quad
\frac{\partial}{\partial y}\qquad\text{oder}\qquad
\frac{d}{dx},\quad\text{und}\quad
\frac{d}{dz},
\]
die die Variable, nach der abgeleitet wird, klar erkennbar macht.
Allerdings ist in der Analysis auch die Notation $f'$ für die
abgeleitete Funktion üblich, aus der man die Variable auch nicht
ablesen kann.

Die spezielle Kennzeichnung des Funktionsargumentes ist aber auch
nicht nötig.
Betrachten wir $x$ als das Funktionsargument, dann ist $f(x)=x$ die
einfachste nichtkonstante Funktion, die wir uns denken können.
Die Ableitung dieser Funktion nach $x$ ist in der klassischen Analysis
ist $f'(x)=1$, wir können daher das Funktionsargument im Funktionenkörper
durch die Relation $Dx=1$ charakterisieren.

\begin{definition}[Differentialkörper]
\index{Differentialkorper@Differentialkörper}%
Ein \emph{Differentialkörper} $K$ ist ein Funktionenkörper mit einer
Derivation.
\end{definition}

Der Hauptsatz der Infinitesimalrechnung besagt, dass das Integral
einer Funktion auch eine Stammfunktion ist.
Die Definition des Integrals ist daher für die Arbeit eines
Computeralgebrasystems nicht relevant, es mur nur eine Stammfunktion
finden können.
Der klassische Begriff der Stammfunktion verwendet die Ableitung.
Für den vorliegenden algebraischen Zusammenhang kann er wie folgt
verallgemeinert werden.

\begin{definition}[Stammfunktion]
Sei $F$ ein Differentialkörper mit Deriviation $D$.
Eine Funktion $\vartheta\in F$ heisst eine \emph{Stammfunktion} von 
$f\in F$, wenn $D\vartheta=f$.
\index{Stammfunktion}%
\end{definition}

%
% Ableitungsregeln
%
\subsubsection{Ableitungsregeln}
Die Produktregel reicht bereits, um die aus dem Analysisunterricht
bekannten Ableitungsregeln zu reproduzieren.

\begin{enumerate}
\item Ableitung einer Potenz: wir betrachten die Potenz $f^n$ eines
Elementes $f\in K$.
Die Derivationseigenschaft liefert
\begin{align*}
D(f^n)
&=
D(\underbrace{f\cdot f\cdot\ldots\cdot f}_{\displaystyle n})
\\
&=
Df\cdot\ldots f
+
f\cdot (Df)\cdot\ldots \cdot f
+
\ldots
+
f\cdot f \cdot\ldots \cdot Df
=
nf^{n-1}\cdot Df.
\end{align*}
Die ist die aus der Analysis wohlbekannte Regel für Ableitungen einer
Potenz.
\item Ableitung der Eins: Die Eins zeichnet sich durch die
Eigenschaft $1^n=1$ aus.
Lässt man die Derivation darauf wirken, erhält man
\[
D(1^n)
=
n1^{n-1}D1
=
D1
\qquad\Rightarrow\qquad
(n-1)D1 = 0.
\]
Da $n-1\ne 1$ ist, folgt $D1=0$, die Ableitung der Eins ist somit $1$.
Da die Derivation $D$ linear ist, gilt für jedes andere Element
$q\in\mathbb{Q}$ sofort $Dq = D(q\cdot 1) = q\cdot D1 = 0$.
\item
Ableitung von $1/f$: Die reziproke Funktion $1/f$ ist charakterisiert
durch die Eigenschaft $f\cdot (1/f)=1$.
Die Ableitung davon ist
\begin{align*}
&&
D(f\cdot(1/f))
&=
(Df)\cdot (1/f) + f\cdot D(1/f)
=
D1
=
0
\\
&\Rightarrow&
f\cdot D(1/f) &= -\frac{Df}{f}
\\
&\Rightarrow&
D(1/f)
&=
-\frac{Df}{f^2},
\end{align*}
dies ist die Ableitungsregel einer Potenz $f^{-1}$.
\item Potenzregel für beliebige Exponenten: Wir wissen bereits,
dass $D(f^n)=nf^{n-1}Df$ für positive $n$.
Für einen negativen Exponenten ist
\[
D(f^{-n})
=
D\biggl(\frac{1}{f^n}\biggr)
=
-
\frac{D(f^n)}{f^{2n}}
=
-\frac{nf^{n-1}Df}{f^{2n}}
=
-
n
f^{-n-1}
Df,
\]
dies ist die erwartete Ableitungsformel für beliebige ganzzahlige 
Exponenten.
\item Ableitungsregel für Polynome:
Sie $p(X) = p(X) = p_0 + p_1X + \dots + p_nX^n \in K[X]$ ein Polynom
mit Koeffizienten im Koeffizientenkörper.
Wir möchten die Ableitung von $Dp(f)$ berechnen.
Nach der Potenzregel gilt
\begin{align*}
Dp(f)
&=
D(p_0+p_1f + p_2f^2 + \ldots + p_{n-1}f^{n-1}+p_nf^n)
\\
&=
p_1Df + p_2\cdot 2f Df + \ldots p_{n-1}(n-1)f^{n-2}Df + p_nnf^{n-1}Df
\\
&=
\bigl(
p_1 + 2p_2 f + 3p_3f^2 + \ldots + p_{n-1}(n-1)f^{n-2} + p_nnf^{n-1}
\bigr)
Df.
\intertext{Das Polynom auf der rechten Seite ist auch bekannt als
das abgeleitete Polynom $p'(X)$, so dass man}
&= p'(f)\, Df
\end{align*}
schreiben kann.
Dies ist formal die Kettenregel für Polynome.
\item Quotientenregel:
Der Quotient $f/g$ kann auch als Produkt $f\cdot (1/g)$ geschrieben
werden.
Daher ist die Ableitung davon
\begin{align*}
D\frac{f}{g}
&=
D\biggl(f\cdot\frac1g\biggr)
=
Df\cdot\frac1g + f\cdot D\frac1g
\\
&=
Df\cdot\frac1g + f\cdot \biggl(-\frac{Dg}{g^2}\biggr)
\\
&=
\frac{Df}{g} - \frac{f\cdot Dg}{g^2}
=
\frac{Df\cdot g - f\cdot Dg}{g^2}.
\end{align*}
Dies ist die bekannte Quotientenregel.
\item Ableitung einer rationalen Funktion:
Eine rationale Funktion ist ein Quotient $p[X]/q[X]$ mit
Polynomen $p,q\in K[X]$.
Wir möchten $p(f)/q(f)$ ableiten:
\begin{align*}
D\biggl(\frac{p(f)}{q(f)}\biggr)
&=
\frac{
(Dp(f))\cdot q(f) - p(f)\cdot Dq(f)
}{
q(f)^2
}
\\
&=
\frac{p'(f)\,Df\cdot q(f)-p(f)\cdot q'(f)\, Df}{q(f)^2}
\\
&=
\frac{p'(f)q(f)-p(f) q'(f)}{q(f)^2}
Df.
\end{align*}
Dies ist die Kettenregel für den Polynombruch $p(X)/q(X)$.
\end{enumerate}

Diese Beispiele zeigen, dass alle Regeln, die man üblicherweise
zur Berechnung von Ableitungen braucht, allein aus der Linearität
und Eigenschaft der Derivation folgen.
Der Funktionenkörper $F$ mit der Derivation $D$ ist daher ein
vollständiges Modell für die algebraischen Eigenschaften der
Ableitung.
Es ist daher auch denkbar, dass das Problem, eine Stammfunktion
zu finden, allein in diesem algebraischen Rahmen gelöst werden
kann.

Der Vollständigkeit halber halten wir hier noch die Definition
des abgeleiteten Polynoms fest.

\begin{definition}[Ableitung eines Polynoms]
\index{Ableitung eines Polynoms}%
Ist $p(X)\in K[X]$ ein Polynom mit Koeffizienten in einem
Koeffizientenkörper mit der Darstellung
\[
p(X)
=
p_0 + p_1X + p_2X^2 + \ldots + p_{n-1}X^{n-1} + p_nX^n,
=
\sum_{k=0}^n p_k X^k
\]
dann heisst
\[
p'(X)
=
p_1 + 2p_2 X + 3p_3 X^2 + \ldots + (n-1)p_{n-1}X^{n-2} + np_n X^{n-1}
=
\sum_{k=1}^n kp_k X^{k-1}
\]
das \emph{abgeleitete Polynom}.
\end{definition}

%
% Ableitung eines algebraischen Elementes
%
\subsubsection{Ableitung eines algebraischen Elementes}
Wenn der Funktionenkörper $F$ um ein algebraisches Element $\alpha$
erweitert wird, dann muss auch die Ableitung $D\alpha$ berechnet
werden können, damit $F(\alpha)$ wieder ein Differentialkörper ist.

Sei als $p(X)\in F[X]$ das Minimalpolynom des Elementes $\alpha$,
es gilt also $p(\alpha)=0$.
Die Ableitung ist nach den Berechungsregeln
\begin{align*}
0
=
Dp(\alpha)
&=
p'(\alpha)D\alpha
+
(Dp_0+Dp_1\alpha + \dots + Dp_n\alpha^n)
\intertext{Da $p$ das Minimalpolynom von $\alpha$ ist, kann $p'(\alpha)$
nicht verschwinden, denn dann wäre $p'$ ein Polynom kleineren Grades,
das auf $\alpha$ verschwindet.
Somit kann man die Gleichung nach $D\alpha$ auflösen und erhält}
D\alpha
&=
-
\frac{Dp_0+(Dp_1)\alpha +\ldots+(Dp_n)\alpha^n}{p'(\alpha)}.
\end{align*}
Die Ableitung eines algebraischen Elementes ist also durch das
Minimalpolynom eindeutig bestimmt.

\begin{beispiel}
Die Quadratwurzel $\!\sqrt{f}$ eines Elementes $f\in F$ hat das
Minimalpolynom $p(X) = X^2 -f \in F[X]$ mit der Ableitung $p'(X)=2X$.
Die Ableitung der Koeffizienten ergibt das Polynom $-Df$.
Damit ist die Ableitung der Quadratwurzel
\begin{align*}
D\!\sqrt{f}
&=
\frac{-Df}{p'(\!\sqrt{f})}
=
-\frac{1}{2\!\sqrt{f}}\cdot Df.
\end{align*}
Dies ist die Kettenregel für die Wurzelfunktion.
\end{beispiel}

\begin{beispiel}
Die $n$-te Wurzel $\sqrt[n]{f}$ eines Elementes $f\in K$ ist algebraisch
mit dem Minimalpolynom $p(X)=X^n -f$ mit Ableitung $p'(X)=nX^{n-1}$ und
dem Polynom mit den abgeleiteten Koeffizienten $-Df$.
Dann ist
\begin{align*}
D\!\sqrt[n]{f}
&=
\frac{-Df}{p'\bigl(\!\sqrt[n]{f}\bigr)}
=
-
\frac{1}{n\bigl(\!\sqrt[n]{f}\bigr)^{n-1}}
\cdot
Df.
\end{align*}
Dies ist die Kettenregel für die $n$-Wurzel.
\end{beispiel}

%
% Konstantenkörper
%
\subsubsection{Konstantenkörper}
Bei der Konstruktion des Funktionskörpers sind wir davon ausgegangen,
dass wir bereits wissen, welche Objekte Konstanten sind.
Die Ableitung $D$ ermöglicht jetzt aber auch, die Konstanten dadurch
zu definieren, dass ihre Ableitung verschwindet.

\begin{definition}[Konstantenkörper]
Sei $F$ ein Differentialkörper, dann heisst
\[
C=\{c\in F\mid Dc=0\}
\]
der Körper der Konstanten oder \emph{Konstantenkörper} in $F$.
\index{Konstantenkorper@Konstantenkörper}%
\end{definition}

Wir müssen uns noch überlegen, dass $C$ wirklich ein Körper ist.
Dazu müssen wir nachrechnen, dass Summe, Produkt und Quotient
konstanter Elemente aus $C$ wieder in $C$ liegt.
Wir wählen daher $a,b\in C$, d.~h.~$Da=0$ und $Db=0$.
Dann gilt auch
\begin{align*}
D(a\pm b)
&=
Da\pm Db = 0\pm 0 = 0
&&\Rightarrow&
a\pm b&\in C
\\
D(ab)
&=
(Da)\cdot b + a\cdot (Db)
=
0\cdot b + a\cdot 0
=
0
&&\Rightarrow&
ab&\in C
\\
D\frac{a}{b}
&=
\frac{(Da)\cdot b - a\cdot (Db)}{b^2}
=
\frac{0\cdot b - a\cdot 0}{b^2}
=
0.
&&\Rightarrow&
\frac{a}{b}&\in C.
\end{align*}
$C$ ist somit ein Körper.
Der Konstantenkörper kann grösser sein als der ursprüngliche
Koeffizientenkörper.

%
% Exponential- und Logarithmusfunktionen
%
\subsection{Exponential- und Logarithmusfunktionen}
Die Exponential- und Logarithmusfunktionen werden in der klassischen
Analysis mithilfe einer Potenzreihe oder durch die Umkehrfunktion
definiert.
Beide Konzepte stehen im algebraischen Umfeld eines Differentialkörpers
nicht zur Verfügung.
Es muss also eine alternative Methode gefunden werden, diese Funktionen
zu charakterisieren.

Die Potenzreihenmethode zur Lösung von Differentialgleichungen
hat in Kapitel~\ref{chapter:differentialgleichungen} gezeigt,
dass die Potenzreihe von $e^x$ eine Folge der Differentialgleichung
für die Exponentialgleichung ist.
Die Differentialgleichung lässt sich auch in einem Differentialkörper
definieren, es liegt daher nahe, diese anstelle der Potenzreihe
zur Definition der Exponentialfunktion zu verwenden.

\begin{definition}[Exponential- und Logarithmusfunktion]
\label{buch:intalg:elementar:def:explog}
Sei $F$ ein Differentialkörper und $f\in F$.
Ein Element $\vartheta\in F$ heisst eine \emph{Exponentialfunktion}
von $f$, wenn $D\vartheta = \vartheta\cdot Df$.
\index{Exponentialfunktion}%
Die Exponentialfunktion wird manchmal auch $\vartheta=e^f$ geschrieben.
$\vartheta$ heisst ein \emph{Logarithmus} von $f$, wenn
$D\vartheta = Df/\vartheta$ oder $\vartheta\cdot D\vartheta = Df$.
\index{Logarithmus}%
Der Logarithmus wird manchmal auch $\log f$ geschrieben.
\end{definition}

%
% Eindeutigkeit
%
\subsubsection{Eindeutigkeit}
Exponential- und Logarithmusfunktion sind durch die Bedingungen der
Definition~\ref{buch:intalg:elementar:def:explog} nicht eindeutig
festgelegt.
Sind $\vartheta_1$ und $\vartheta_2$ zwei Logarithmen von $f\in F$,
dann ist
\[
D(\vartheta_1-\vartheta_2)
=
\frac{Df}{f}-\frac{Df}{f}
=
0,
\]
die Differenz ist also eine Konstante $\vartheta_1-\vartheta_2\in C$.
Eine Logarithmusfunktion ist also nur bis auf eine Konstante in $C$
bestimmt.
Verwendet man die Notation $\log f$ suggeriert das zwar eine eindeutig
bestimmte Funktion, es kann aber nur eine Menge von Funktionen
$\vartheta_1 + C$ gemeint sein.

Sind andererseits $\varphi_1$ und $\varphi_2$ Exponentialfunktionen
von $f$, dann ist
\[
D
\frac{\varphi_1}{\varphi_2}
=
\frac{D\varphi_1\cdot\varphi_2 - \varphi_1\cdot D\varphi_2}{\varphi_2^2}
=
\frac{
\varphi_1 Df\varphi_2 - \varphi_1 \varphi_2 Df
}{\varphi_2^2}
=
0,
\]
der Quotient $\varphi_1/\varphi_2$ ist also eine Konstante.
Eine Exponentialfunktion ist also nur bis auf einen konstanten
Faktor in $C$ bestimmt.
Die Notation $e^f$ suggeriert eine eineutig bestimmte Funktion,
gemeint sein kann aber nur eine Menge $C\varphi_1$ von Funktionen in $F$.

%
% Logarithmus als Umkehrfunktion der Exponentialfunktion
%
\subsubsection{Logarithmus als Umkehrfunktion der Exponentialfunktion}
Das Konzept der Zusammensetzung von Funktionen kommt in einem
Differentialkörper nicht vor, so dass die Frage, ob der
Logarithmus die Umkehrfunktion der Exponentialfunktion ist,
nicht sinnvoll gestellt werden kann.
Der Begriff der Umkehrfunktion ist nur sinnvoll, wenn man
\index{Umkehrfunktion}%
Funktionen zusammensetzen kann.

Trotzdem kann man eine Interpretation finden, nach der der Logarithmus
eine Art Umkehrfunktion ist.
Da die Exponentialfunktion und der Logarithmus durch die
Definition über ihre differentiellen Eigenschaften nicht eindeutig 
bestimmt sind, kann auch die Umkehrbeziehung nur im Sinne einer
Differentialbeziehung gelten.

Sei jetzt $f\in F$ und $\vartheta$ eine Exponentialfunktion von $f$.
Nach Definition erfüllt sie $D\vartheta=\vartheta\cdot Df$.
Weiter sei $\varphi$ ein Logarithmus von $\vartheta$, d.~h.~nach
Definition ist $D\varphi=D\vartheta/\vartheta$.
Setzt man die definierende Eigenschaft der Exponentialfunktion
in die Ableitung des Logarithmus ein, erhält man
\[
D\varphi
=
\frac{D\vartheta}{\vartheta}
=
\frac{\vartheta\cdot Df}{\vartheta}
=
Df.
\]
Dies bedeutet, dass $\varphi$ bis auf eine additive Konstante mit
$f$ übereinstimmt.
Mit der konventionellen Schreibweise für Exponentialfunktion und
Logarithmus kann man dies auch als
\[
D\log e^f
=
\frac{De^f}{e^f}
=
\frac{e^f Df}{e^f}
=
Df
\]
schreiben.
Dies bedeutet, dass $\log e^f = f+c$ mit einer Konstante $c\in C$ ist.

Umgekehrt ist
\begin{align*}
D\frac{e^{\log f}}{f}
&=
\frac{(De^{\log f})\cdot f - e^{\log f}\cdot Df}{f^2}
\\
&=
\frac{e^{\log f}D(\log f)\cdot f - e^{\log f}\cdot Df }{f^2}
\\
&=
\frac{e^{\log f} (Df/f)\cdot f - e^{\log f}\cdot Df }{f^2}
=
0,
\end{align*}
$e^{\log f}$ ist also bis auf einen konstanten Faktor die ursprüngliche
Funktion $f$.

%
% Logarithmus und Exponentialgesetze
%
\subsubsection{Logarithmus- und Exponentialgesetze}
Für die algebraische Rechnung sind Regeln nötig, mit denen Ausdrücke
mit Logarithmen und Exponentialfunktionen umgeformt werden können.
Dazu gehören die Logarithmus- und Exponentialgesetze
\index{Logarithmusgesetz}%
\index{Exponentialgesetz}%
\[
\log ab = \log a + \log b
\qquad\text{und}\qquad
e^{a+b}=e^ae^b.
\]
Da weder Logarithmus noch Exponentialfunktion eindeutig festgelegt
sind, kann man so eine einfache Beziehung nicht mehr erwarten.

Sind $f_1,f_2\in F$ Funktionen in $F$ und $\vartheta_1$ und $\vartheta_2$
Exponentialfunktionen von $f_1$ bzw.~$f_2$, dann gilt
\[
D(\vartheta_1\vartheta_2)
=
(D\vartheta_1)\vartheta_2
+
\vartheta_1(D\vartheta_1)
=
\vartheta_1 (Df_1) \vartheta_2
+
\vartheta_1\vartheta_2 (Df_2)
=
\vartheta_1\vartheta_2 D(f_1+f_2).
\]
Das Produkt $\vartheta = \vartheta_1\vartheta_2$ ist daher wegen
\[
D\vartheta = \vartheta D(f_1+f_2)
\]
eine Exponentialfunktion von $f_1+f_2$.
In der konventionellen Notation wird dies als
\[
e^{f_1} e^{f_2}
=
e^{f_1+f_2}
\]
geschrieben.
Das Produkt der Exponentialfunktionen der Faktoren ist also eine
Exponentialfunktion der Summe.
Das Exponentialgesetzt ist also in diesem erweiterten Sinn immer noch
gültig.

Seien jetzt umgekehrt $\vartheta_1$ und $\vartheta_2$ Logarithmen von
$f_1$ bzw.~$f_2$.
Dann gilt
\[
D(\vartheta_1+\vartheta_2)
=
\frac{Df_1}{f_1}
+
\frac{Df_2}{f_2}
=
\frac{Df_1\cdot f_2+f_1\cdot Df_2}{f_1f_2}
=
\frac{D(f_1f_2)}{f_1f_2}.
\]
Schreibt man $f=f_1f_2$, dann ist die Summe
$\vartheta=\vartheta_1+\vartheta_2$ wegen
\[
D\vartheta = \frac{Df}{f}
\]
ein Logarithmus des Produktes $f$.
In der konventionellen Notation bedeutet dies, dass
\[
\log (f_1f_2)
=
\log f_1 + \log f_2
\]
ist.
Die Summe der Logarithmen der Faktoren ist also ein Logarithmus
des Produktes.


%
% Exponentialdarstellung der trigonometrischen Funktionen
%
\subsection{Exponentialdarstellung der trigonometrischen Funktionen
\label{buch:intalg:elementar:subsection:trigo}}
Die trigonometrischen Funktionen wurden ursprünglich für die Zwecke
der Beschreibung der Himmelskugel entwickelt und in der ebenen
Geometrie bei der Berechnung von Dreiecken angewendet.
Mit der eulerschen Formel $e^{it}=\cos t + i\sin t$ wird jedoch 
\index{eulersche Formel}%
auch eine Beziehung zur komplexen Exponentialfunktion offengelegt.
Sie ermöglicht die trigonometrischen Funktionen und ihre Umkehrfunktionen
ausschliesslich die Exponentialfunktion und ihre Umkehrfunktion, den
Logarithmus, zu beschreiben.

%
% Hyperbolische Funktionen
%
\subsubsection{Hyperbolische Funktionen}
Bei den hyperbolischen Funktionen ist die Situation viel einfacher,
da diese Funktionen bereits durch die reelle Exponentialfunktion
definiert sind.

\begin{definition}[Hyperbolische Funktionen]
\label{buch:intalg:elementar:trigo:def:hyperbolisch}
Die \emph{hyperbolischen Funktionen} sind durch
\index{hyperbolische Funktion}%
\begin{align*}
\sinh x &= \frac{e^x-e^{-x}}2,
&
\cosh x &= \frac{e^x+e^{-x}}2,
&
\tanh x &= \frac{e^x-e^{-x}}{e^x+e^{-x}},
&
\coth x &= \frac{e^x+e^{-x}}{e^x-e^{-x}}
\end{align*}
definiert für beliebige Argumente $x\in \mathbb{C}$.
\index{sinh}%
\index{cosh}%
\index{tanh}%
\index{coth}%
\end{definition}

Die hyperbolischen Funktionen sind besonderes nützlich für die Lösung
von Differentialgleichungen zweiter Ordnung, weil sie ähnlich einfache 
Anfangsbedingungen wie die Funktionen $\cos x$ und $\sin x$ haben.
Die Differentialgleichung $y''=y$ hat die Exponentialfunktionen 
$y_+(x)=e^x$ und $y_-(x)=e^{-x}$ als linear unabhängige Lösungen.
Um jedoch eine Lösung der Differentialgleichung zu den Anfangsbedingungen
$y(0)=y_0$ und $y'(0)=y_1$ zu finden, muss man eine Linearkombination
\[
y(x) = A y_+(x) + B y_-(x)
\]
finden, die die korrekten Anfangsbedingungen erfüllt. 
Dazu setzt  man die Linearkombination in die Anfangsbedinungen ein
und erhält das Gleichungssystem
\[
\renewcommand{\arraycolsep}{2pt}
\left.
\begin{array}{rcrcrcr}
y(0)  &=& A y_+(0)  &+& B y_-(0)  &=& y_0 \\
y'(0) &=& A y_+'(0) &+& B y_-'(0) &=& y_1 \\
\end{array}
\quad
\right\}
\qquad
\Rightarrow
\qquad
\left\{
\quad
\begin{array}{rcrcl}
A &+& B &=& y_0 \\
A &-& B &=& y_1.
\end{array}
\right.
\]
Durch Addition bzw.~Subtraktion der beiden Gleichungen ist das 
Gleichungssystem einfach zu lösen, man findet die Koeffizienten
\[
A
=
\frac{y_0+y_1}2
\qquad\text{und}\qquad
B
=
\frac{y_0-y_1}2
\]
und daraus die Lösung
\begin{align*}
y(x)
=
Ay_+(x) + By_-(x)
&=
\frac{y_0+y_1}2 e^x
+
\frac{y_0-y_1}2 e^{-x}
\\
&=
y_0
\frac{e^x+e^{-x}}2
+
y_1
\frac{e^x-e^{-x}}2
=
y_0\cosh x + y_1 \sinh x
\end{align*}
der Differentialgleichung.
Die hyperbolischen Funktionen tauchen also genauso natürlich in der
Lösung der Differentialgleichung $y''=y$ auf, wie die trigonometrischen
Funktionen in der Lösung der Schwingungsdifferentialgleichung $y''=-y$
auftreten.
Der Grund dafür ist natürlich, dass die Funktionen die gleichen
Anfangsbedingungen
\[
\begin{aligned}
\cos  0 &= 1,& \cos'  0 &= 0,\\
\sin  0 &= 0,& \sin'  0 &= 1
\end{aligned}
\qquad\text{bzw.}\qquad
\begin{aligned}
\cosh 0 &= 1,& \cosh' 0 &= 0,\\
\sinh 0 &= 0,& \sinh' 0 &= 1
\end{aligned}
\]
haben.

%
% Die hyperbolischen Umkehrfunktionen
%
\subsubsection{Die hyperbolischen Umkehrfunktionen}
Die Logarithmusfunktion ist die Umkehrfunktion der Exponentialfunktion.
Mit ihr lassen sich auch die hyperbolischen Funktionen invertieren.

Um die Umkehrfunktion $\operatorname{arcosh}y$ zu finden, muss die
reelle Zahl $x$ gefunden werden, die $\cosh x = y$ gilt.
Aus der Definition folgt die Gleichung
\begin{align*}
\frac{e^x + e^{-x}}2 &= y.
\intertext{Da die Exponentialfunktion nicht verschwinden kann, verlieren
wir nichts, wenn wir mit $e^x$ multiplizieren. 
Dies ergibt}
\frac{(e^x)^2 + 1}2 &= y e^x
\intertext{oder als quadratische Gleichung für $e^x$ geschrieben:}
(e^x)^2 - 2y e^x + 1 &= 0.
\end{align*}
Die Lösung der quadratischen Gleichung ist
\begin{equation}
e^x
=
y\pm \!\sqrt{\smash{y^2}-1\raisebox{1pt}{\mathstrut}}
\qquad
\Rightarrow
\qquad
x
=
\log\Bigl(
y\pm \!\sqrt{\smash{y^2}-1\raisebox{1pt}{\mathstrut}}
\Bigr).
\label{buch:intalg:elementar:eqn:arcosh}
\end{equation}
Da $\cosh x$ eine gerade Funktion ist, gibt es zu den meisten
$y$-Werten zwei $x$-Werte.
Da $\cosh x$ eine gerade Funktion ist, würden wir erwarten,
dass die beiden Werte in \eqref{buch:intalg:elementar:eqn:arcosh}
entgegengesetzt sind.
Dazu müssen aber die Argumente des Logarithmus reziprok sein.
Tatsächlich
\begin{align*}
\frac{1}{y\pm\!\sqrt{\smash{y^2}-1\raisebox{1.5pt}{\mathstrut}}}
&=
\frac{
y\mp\!\sqrt{\smash{y^2}-1\raisebox{1.5pt}{\mathstrut}}
}{
(
y\pm\!\sqrt{\smash{y^2}-1\raisebox{1.5pt}{\mathstrut}}
)(
y\mp\!\sqrt{\smash{y^2}-1\raisebox{1.5pt}{\mathstrut}}
)
}
\\
&=
\frac{
y\mp\!\sqrt{\smash{y^2}-1\raisebox{1.5pt}{\mathstrut}}
}{
y^2-(y^2-1)
}
\\
&=
y\mp\!\sqrt{\smash{y^2}-1\raisebox{1.5pt}{\mathstrut}}.
\end{align*}
Die beiden Argumente des Logarithmus sind also tatsächlich
reziprok.

Auch für die hyperbolische Sinusfunktion kann mit der gleichen
Methode ein Ausdruck für die Umkehrfunktion gefunden werden.
Aus
\[
\sinh x = y
\qquad\Rightarrow\qquad
\frac{e^x - e^{-x}}2 = y
\]
folgt durch Multiplikation mit $2e^x$
\begin{align*}
(e^x)^2 -2ye^x - 1 &= 0
\notag
\intertext{mit der Lösung}
e^x &= y\pm\!\sqrt{\smash{y^2}+1\raisebox{1.5pt}{\mathstrut}}.
\label{buch:intalg:elementar:trig:eqn:arsinh}
\end{align*}
Da das negative Zeichen auf der rechten Seite zu einem negativen Wert 
führt, bleibt nur die Lösung
\[
x
=
\log\Bigl(y+\!\sqrt{\smash{y^2}+1\raisebox{1.5pt}{\mathstrut}}\Bigr)
.
\]

Für die hyperbolische Tangens- und Kotangensfunktion ist der
entsprechende Ansatz
\begin{align*}
\tanh x
=
\frac{e^x-e^{-x}}{e^x+e^{-x}}
&=
y
&
&\text{bzw.}&
\coth x
=
\frac{e^x+e^{-x}}{e^x-e^{-x}}
&=
y
\intertext{oder nach Multiplizieren mit dem Nenner}
e^x-e^{-x}
&=
ye^x+ye^{-x}
&
&&
e^x+e^{-x}
&=
ye^x-ye^{-x}
\intertext{und $e^x$}
(e^x)^2-1
&=
y(e^x)^2+y
&
&&
(e^x)^2+1
&=
y(e^x)^2-y.
\intertext{Umgeschrieben als quadratische Gleichung für $e^x$ ist dies}
(1-y) (e^x)^2
&=
1 +y
&
&&
(1-y) (e^x)^2
&=
-1-y.
\intertext{Mit der Quadratwurzel findet man}
e^x
&=
\!\sqrt{\frac{1+y}{1-y}}
&
&\text{bzw.}&
e^x
&=
\!\sqrt{\frac{y+1}{y-1}}
\intertext{und schliesslich mit dem Logarithmus den Wert}
x
=
\operatorname{artanh}y
&=
\log
\!\sqrt{\frac{1+y}{1-y}}
&
&\text{bzw.}&
x
=
\operatorname{arcoth}y
&=
\log
\!\sqrt{\frac{y+1}{y-1}}
\\
&=
\frac12\log
\frac{1+y}{1-y}
&
&&
&=
\frac12\log
\frac{y+1}{y-1}
\end{align*}
der Umkehrfunktion der hyperbolischen Tangens- bzw.~Kotangensfunktion.
\index{artanh}%
\index{arcoth}%

\begin{satz}
\label{buch:intalg:elementar:trigo:satz:invhyperbolisch}
Die Umkehrfunktionen der hyperbolischen Funktionen sind
\begin{align*}
\operatorname{arsinh}y
&=
\log\Bigl(
y + \!\sqrt{\smash{y^2}+1\raisebox{1.5pt}{\mathstrut}}
\Bigr),
&
\operatorname{arcosh}y
&=
\log\Bigl(
y + \!\sqrt{\smash{y^2}-1\raisebox{1pt}{\mathstrut}}
\Bigr),
\\
\operatorname{artanh}y
&=
\frac12\log
\frac{1+y}{1-y},
&
\operatorname{arcoth}y
&=
\frac12\log
\frac{y+1}{y-1}.
\end{align*}
\end{satz}

Der Satz zeigt also, dass nicht nur die hyperbolischen Funktionen
sondern auch ihre Umkehrfunktionen ausschliesslich durch algebraische
Funktionen und den Logarithmus definiert werden können.

Die Funktionen $\operatorname{tanh}x$ und $\operatorname{coth}x$ sind
reziprok.
Die beiden Funktionen $\operatorname{artanh}y$ und
$\operatorname{arcoth}y$
müssen also unter der Substitution $y\to 1/y$ ineinander übergehen.
Tatsächlich wird das Argument des Logarithmus in der Areatangensfunktion
unter der Substitution zu
\[
\frac{1+y}{1-y}
\to
\frac{1+1/y}{1-1/y}
=
\frac{y+1}{y-1},
\]
dem Argument des Logarithmus der Areakotangensfunktion.

%
% Sinus und Kosinus
%
\subsubsection{Sinus und Kosinus}
Aus der Identität
\[
e^{ix}
=
\cos x + i \sin x
\]
folgt zunächst mittels der Substitution $x\to -x$
\[
e^{-ix}
=
\cos x - i \sin x.
\]
Diese beiden Formeln bilden ein lineares Gleichungssystem
\[
\renewcommand{\arraycolsep}{2pt}
\begin{array}{rcrcl}
\cos x &+& i\sin x &=& e^{\phantom{-}ix} \\
\cos x &-& i\sin x &=& e^{-ix}
\end{array}
\]
für die Funktionen $\cos x$ und $\sin x$.
\index{sin}%
\index{cos}%
Es kann durch Addition bzw.~Subtraktion gelöst werden.
Die Lösungen sind
\[
\cos x = \frac{e^{ix}+e^{-ix}}{2}
\qquad\text{und}\qquad
\sin x = \frac{e^{ix}-e^{-ix}}{2i}.
\]
Die Funktionen können also mit der Exponentialfunktion in der gleichen
Form ausgedrückt werden, in der die hyperbolischen Funktionen 
definiert sind.
Durch Vergleich mit den Formeln in der
Definition~\ref{buch:intalg:elementar:trigo:def:hyperbolisch}
ergeben sich die Beziehungen
\[
\cos x = \cosh ix
\qquad\text{und}\qquad
\sin x = \frac{1}{i} \sinh ix.
\]
Die Multiplikation mit $i$ entspricht geometrisch einer Drehung
der komplexen Zahlenebene um $90^\circ$.
Man kann also sagen, dass trigonometrischen Funktionen durch Drehung
der Argumentebene um $90^\circ$ aus den hyperbolischen Funktionen
hervorgehen.

Für die Umkehrfunktionen $\arccos y$ und $\arcsin y$ sind die gleichen
algebraischen Umformungen anwendbar, die auf die Umkehrfunktionen
$\operatorname{arcosh}y$ und $\operatorname{arsinh}y$ der hyperbolischen
Funktionen geführt haben.
\index{arccos}%
\index{arcsin}%
Aus
\input{chapters/070-intalg/fig/fig-invcis.tex}%
\begin{align}
\cos x = \frac{e^{ix}+e^{-ix}}{2} &= y
&
&\text{bzw.}&
\sin x = \frac{e^{ix}-e^{-ix}}{2i} &= y
\notag
\intertext{folgen durch Multiplikation mit dem Nenner und mit $e^{ix}$
die quadratischen Gleichungen}
(e^{ix})^2 -2ye^{ix} + 1 &= 0
&
&&
(e^{ix})^2 - 2iy e^{ix} - 1 &= 0
\notag
\intertext{für $e^{ix}$ mit den Lösungen}
e^{ix}
&=
y \pm \!\sqrt{\smash{y^2}-1\raisebox{1.5pt}{\mathstrut}}
&
&&
e^{ix}
&=
iy\pm \!\sqrt{-\smash{y^2}+1\raisebox{1.5pt}{\mathstrut}}.
\notag
\intertext{Da die Kosinuswerte Betrag nicht grösser als $1$ haben
können, ist der Radikand in der linken Gleichung immer negativ.
Wir können das Vorzeichen kehren, indem wir den Faktor $i$ vor die
Wurzel ziehen:}
e^{ix}
&=
y \pm i\!\sqrt{1-\smash{y^2}\raisebox{1.5pt}{\mathstrut}}
&
&&
e^{ix}
&=
iy\pm \!\sqrt{1-\smash{y^2}\raisebox{1.5pt}{\mathstrut}}.
\label{buch:intalg:elementar:trigo:eqn:invcisrhs}
\end{align}
Abbildung~\ref{buch:intalg:elementar:trigo:fig:invcis} zeigt 
in rot die rechten Seiten der Gleichungen
\eqref{buch:intalg:elementar:trigo:eqn:invcisrhs}.
Da die Punkte auf dem Einheitskreis liegen, ist der Logarithmus
der in blau eingezeichnete Winkel $x$.

\begin{satz}
\label{buch:intalg:elementar:trigo:satz:trigonometrisch}
Die trigonometrischen Funktionen
\begin{align*}
\cos x &= \frac{e^{ix}+e^{-ix}}{2}
&&\text{und}&
\sin x &= \frac{e^{ix}-e^{-ix}}{2i}
\intertext{in Exponentialdarstellung haben die logarithmischen
Umkehrfunktionen}
\arccos y
&=
\frac{1}{i}
\log
\Bigl(
y + i\!\sqrt{1-\smash{y^2}\raisebox{1.5pt}{\mathstrut}}
\Bigr)
&&\text{bzw.}&
\arcsin y
&=
\frac{1}{i}
\log
\Bigl(
iy + \!\sqrt{1-\smash{y^2}\raisebox{1.5pt}{\mathstrut}}
\Bigr).
\end{align*}
\end{satz}

%
% Tangens und Kotangens
%
\subsubsection{Tangens und Kotangens}
Tangens und Kotangens sind Quotienten der Sinus- und Kosinusfunktionen
\input{chapters/070-intalg/fig/fig-invtan.tex}%
oder in Exponentialdarstellung
\begin{align*}
\tan x
&=
\frac1{i}\frac{e^{ix}-e^{-ix}}{e^{ix}+e^{-ix}}
&
&\text{bzw.}&
\cot x
&=
i\frac{e^{ix}+e^{-ix}}{e^{ix}-e^{-ix}}
\end{align*}
Für die Umkehrfunktionen setzt man wieder 
\index{arctan}%
\begin{align*}
\tan x
= 
\frac1i\frac{e^{ix}-e^{-ix}}{e^{ix}+e^{-ix}}
&=
y
&&&
\cot x
=
i\frac{e^{ix}+e^{-ix}}{e^{ix}-e^{-ix}}
&=
y
\intertext{und multipliziert mit dem Nenner}
e^{ix}-e^{-ix}
&=
iy
(e^{ix}+e^{-ix})
&&&
e^{ix}+e^{-ix}
&=
-iy
(e^{ix}-e^{-ix})
\\
(1-iy)e^{ix}
&=
(1+iy)e^{-ix}
&&&
(iy-1)e^{ix}
&=
-(1+iy)e^{-ix}
\intertext{Durch Multiplizieren mit $e^{ix}$ kann man nach $e^{2ix}$
auflösen und findet}
e^{2ix}
&=
\frac{1+iy}{1-iy}
&&&
e^{2ix}
&=
\frac{1-iy}{1+iy}
\\
x
&=
\frac1{2i}
\log
\frac{1+iy}{1-iy}
&&\text{bzw.}&
x
&=
\frac1{2i}
\log
\frac{1-iy}{1+iy}.
\end{align*}
In Abbildung~\ref{buch:intalg:elementar:trigo:fig:invtan} sind
die Zahlen $1+iy$ und $1-iy$ dargestellt.
Der gesuchte Winkel $x$ tritt beim Ursprung auf.
Da die Multiplikation mit einer komplexen Zahl eine Drehungstreckung ist,
muss der Quotient den Drehstreckungsfaktor zurückliefen.
Der Quotient von $1+iy$ und $1-iy$ ist ein komplexe Zahl vom Betrag $1$,
die Drehung von $1-iy$ auf $i+iy$ beschreibt.
Der zugehörige Drehwinkel ist $2x$.
Somit ist der gesuchte Winkel tatsächlich der halbe Logarithmus
des Quotienten der beiden Zahlen.

\begin{satz}
Die trigonometrischen Funktionen
\begin{align*}
\tan x
&=
\frac1{i}\frac{e^{ix}-e^{-ix}}{e^{ix}+e^{-ix}}
&
&\text{und}&
\cot x
&=
i\frac{e^{ix}+e^{-ix}}{e^{ix}-e^{-ix}}
\intertext{in Exponentialdarstellung haben die logarithmische Umkehrfunktion}
\arctan y
&=
\frac1{2i}
\log
\frac{1+iy}{1-iy}
&&\text{bzw.}&
\operatorname{arccot} y
&=
\frac1{2i}
\log
\frac{1-iy}{1+iy}.
\end{align*}
\end{satz}

%
% Stammfunktion einer rationalen Funktion
%
\subsubsection{Stammfunktion einer rationalen Funktion}
In Abschnitt~\ref{buch:integration:residuen:section:beispiel}
wurde die Aufgabe gestellt, das die Stammfunktion
\[
\int \frac{dx}{x^2+1}
\]
zu bestimmen.
Die Partialbruchzerlegung des Integranden ist
\[
\frac{1}{x^2+1}
=
\frac1{2i}
\biggl(
\frac{1}{x-i}
-
\frac{1}{x+i}
\biggr).
\]
In dieser Form ist die Stammfunktion einfach zu bestimmen, sie ist
\begin{align*}
\int
\frac{1}{x^2+1}
\,dx
&=
\frac{1}{2i}
\bigl(
\log(x-i)
-
\log(x+i)
\bigr)
=
\frac{1}{2i}
\log
\frac{x-i}{x+i}
=
\frac{1}{2i}
\log
\frac{1+ix}{-1+ix}.
\intertext{%
In dieser Form stimmt das Resultat nicht mit der Form des erwarteten
Resultates, nämlich der Arkustangens-Funktion überein.
Der Nenner hat das falsche Vorzeichen.
Multiplikation mit $-1$ im Argument des Logarithmus ändert aber den
Wert der Logarithmusfunktion nur um die Konstante $\pi$, also kann
man immer noch schliessen}
&=
\frac{1}{2i}
\log
\frac{1+ix}{1-ix}
+C
=
\arctan x + C.
\end{align*}
Das Beispiel illustriert, dass sich Integrale von rationalen Funktionen
über $\mathbb{R}$ tatsächlich nur mit der komplexen Logarithmus-Funktion
berechnen lassen.

%
% Die Differentialkörper der hyperbolischen Funktionen
%
\subsubsection{Die Differentialkörper der hyperbolischen Funktionen}
Der Körper $\mathbb{Q}(x)$ enthält weder die trigonometrischen noch
die hyperbolischen Funktionen.
Wie früher betrachten wir die Derivation $D$ mit $Dx=1$, die die
Ableitung nach $x$ symbolisiert.

Um die hyperbolischen Funktionen zu $\mathbb{Q}(x)$ hinzuzufügen,
muss ein Exponentialfunktion $e(x)$ hinzugefügt werden.
Damit lassen sich dann die Funktionen
\[
\cosh x
=
\frac12
\biggl(
e(x) + \frac{1}{e(x)}
\biggr),
\quad
\sinh x
=
\frac12
\biggl(
e(x) - \frac{1}{e(x)}
\biggr),
\quad
\tanh x
=
\frac{
\displaystyle
e(x)-\frac{1}{e(x)}
}{
\displaystyle
e(x)+\frac{1}{e(x)}
}
=
\frac{e(x)^2-1}{e(x)^2+1}
.
\]
Die bekannte Relation
\[
(\cosh x)^2
-
(\sinh x)^2
=
\frac14\biggl(
e(x)^2 + 2e(x)\frac{1}{e(x)} + \frac{1}{(e(x))^2}
\biggr)
-
\frac14\biggl(
e(x)^2 - 2e(x)\frac{1}{e(x)} + \frac{1}{(e(x))^2}
\biggr)
=
1
\]
folgt allein durch algebraische Umformung.
Die Ableitungseigenschaften der Funktion $e(x)$ im Differentialkörper
als Exponentialfunktion von $x$ reichen aus, um die Ableitungsregeln
\begin{align*}
D\cosh x
&=
\frac12\biggl(
De(x) + D\frac1{e(x)}
\biggr)
=
\frac12
\biggl(
e(x) - \frac{1}{e(x)}
\biggr)
=
\sinh x,
\\
D\sinh x
&=
\frac12
\biggl(
De(x) - D\frac{1}{e(x)}
\biggr)
=
\frac12
\biggl(
e(x) + \frac{1}{e(x)}
\biggr)
=
\cosh x
\end{align*}
für die hyperbolischen Funktionen zu verifizieren.
Auch die Ableitungsregel für den hyperbolischen Tangens können wir
leicht nachrechnen:
\begin{align*}
D\tanh x
&=
D
\frac{\sinh x}{\cosh x}
\\
&=
\frac{(D\sinh x)\cdot \cosh x - \sinh x\cdot(D\cosh x)}{(\cosh x)^2}
\\
&=
\frac{(\cosh x)^2 - (\sinh x)^2}{(\cosh x)^2}
\\
&=
\frac{1}{(\cosh x)^2}.
\end{align*}
Dies ist die bekannte Ableitung der hyperbolischen Tangensfunktion.
Diese Beispiele zeigen, dass alle Eigenschaften der hyperbolischen
Funktionen von den Details der Funktion $e(x)$ unabhängig sind.

Für die hyperbolischen Umkehrfunktionen müssen gemäss den Formeln
von Satz~\ref{buch:intalg:elementar:trigo:satz:invhyperbolisch}
zunächst die Wurzeln
\[
s_+(x) = \sqrt{\smash{y^2} + 1\raisebox{1.5pt}{\mathstrut}}
\qquad\text{und}\qquad
s_-(x) = \sqrt{\smash{y^2} - 1\raisebox{1.5pt}{\mathstrut}}
\]
hinzufügen, mit denen anschliessend im Funktionenkörper
$\mathbb{Q}(x,e(x),s_+(x),s_-(x))$
die Funktionen
\[
x + s_+(x)
\qquad\text{und}\qquad
x + s_-(x)
\]
gebildet werden können.
Indem wir die Logarithmen $\operatorname{arsinh}(x) = \log(x+s_+(x))$
und $\operatorname{arcosh}(x) = \log(x+s_-(x))$
hinzufügen, erhalten wir den Differentialkörper
\[
\mathbb{Q}(x,e(x),s_\pm(x),
\operatorname{arsinh}(x),\operatorname{arcosh}(x)),
\]
die die gesuchten Umkehrfunktionen entält.
Die hyperbolischen Funktionen und ihre Umkehrfunktionen können also
allein durch den Prozess der algebraischen, exponentiellen oder
logarithmischen Erweiterung erhalten werden.

%
% Die Differentialkörper der trigonometrischen Funktionen
%
\subsubsection{Die Differentialkörper der trigonometrischen Funktionen}
Ein Differentialkörper für die trigonometrischen Funktionen und ihre
Umkehrfunktionen lässt sich ganz analog zur Konstruktion für die 
hyperbolischen Funktionen im vorangegangenen Abschnitt konstruieren.
Damit die Konstruktion möglich wird, muss dem Körper $\mathbb{Q}(x)$
im ersten Schritt die neue Konstante $i$ durch algebraische
Erweiterung mit dem Minimalpolynom $x^2+1$ hinzugefügt werden.
Der Ausgangskörper für die Konstruktion ist also $\mathbb{Q}(x,i)$.

Statt der reellen Exponentialfunktion fügen wir im nächsten Schritt 
dem Körper die Exponentialfunktion $e(x)$ von $ix$ hinzu.
Sie erfüllt die definierende Eigenschaft $De(x) = e(x)\cdot D(ix) = ie(x)$.
Damit können dann die trigonometrischen Funktionen
\begin{align*}
\sin x
&=
\frac12\biggl(
e(x)+\frac{1}{e(x)}
\biggr),
&
\cos x
&=
\frac1{2i}\biggl(
e(x)-\frac{1}{e(x)}
\biggr),
&
\tan x
&=
\frac{1}{i}
\frac{
\displaystyle e(x)-\frac{1}{e(x)}
}{
\displaystyle e(x)+\frac{1}{e(x)}
}
=
\frac{1}{i}\frac{e(x)^2-1}{e(x)^2+1}
\end{align*}
durch algebraische Operationen in diesem Körper konstruiert werden.
Der ``Satz des Pythagoras'' ist eine rein algebraische Eigenschaft,
es gilt nämlich
\begin{align*}
(\sin x)^2 + (\cos x)^2
&=
\frac1{4i^2}
\biggl(
e(x)^2 - 2 + \frac{1}{e(x)^2}
\biggr)
+
\frac14
\biggl(
e(x)^2 + 2 + \frac{1}{e(x)^2}
\biggr)
\\
&=
-
\frac1{4}
\biggl(
e(x)^2 - 2 + \frac{1}{e(x)^2}
\biggr)
+
\frac14
\biggl(
e(x)^2 + 2 + \frac{1}{e(x)^2}
\biggr)
=
1.
\end{align*}
Auch die Ableitungsregeln folgen direkt aus der Definition
\begin{align*}
D\sin x
&=
\frac{1}{2i}
\biggl(
De(x) - D\frac{1}{e(x)}
\biggr)
=
\frac{1}{2i}
\biggl(
ie(x) + i\frac{1}{e(x)}
\biggr)
=
\cos x
\\
D\cos x
&=
\frac12
\biggl(
De(x) + D\frac{1}{e(x)}
\biggr)
=
\frac12
\biggl(
ie(x) - i\frac{1}{e(x)}
\biggr)
=
-
\frac1{2i}
\biggl(
e(x) - \frac{1}{e(x)}
\biggr)
=
-
\sin x.
\end{align*}
Die Funktionen $\sin x$ und $\cos x$ sind also in der Körpererweiterung
$\mathbb{Q}(x,i,e(x))$ enthalten.

Für die Umkehrfunktionen sind nach
Satz~\ref{buch:intalg:elementar:trigo:satz:trigonometrisch}
zusätzlich das algebraische Element
\[
s(x)
=
\! \sqrt{1-\smash{x^2}\raisebox{1.5pt}{\mathstrut}}
\]
hinzufügen.
Der Differentialkörper $\mathbb{Q}(x,i,s(x))$ enthält dann auch
die Funktionen
\[
x+is(x)
=
x + i\!\sqrt{1-\smash{y^2}\raisebox{1.5pt}{\mathstrut}}
\qquad
\text{und}
\qquad
ix+s(x)
=
ix + \!\sqrt{1-\smash{y^2}\raisebox{1.5pt}{\mathstrut}},
\]
da sie durch algebraische Operationen entstehen.
Schliesslich müssen nur noch die Logarithmen
\[
\vartheta_1(x)
=
\log(y+is(x))
\qquad
\text{und}
\qquad
\vartheta_2(x)
=
\log(iy+s(x))
\]
dieser Funktionen hinzugefügt werden, damit dann in
$\mathbb{Q}(x,i,s(x),\vartheta_1(x),\vartheta_2(x))$
die Funktionen
\[
\arccos x = \frac{1}{i}\vartheta_1(x)
\qquad
\text{und}
\qquad
\arcsin x = \frac{1}{i}\vartheta_2(x)
\]
enthält.

Wie bei den hyperbolischen Funktionen können wir schliessen, dass
den Differentialkörper $\mathbb{Q}(x,i,s(x),\vartheta_1(x),\vartheta_2(x))$
die inversen trigonometrischen Funktionen entält.
Die trigonometrischen Funktionen ebenso wie ihre Umkehrfunktionen
sind daher elementare Funktionen.

%
% Eindeutigkeit der trigonometrischen und hyperbolischen Funktionen
%
\subsubsection{Eindeutigkeit der trigonometrischen und hyperbolischen
Funktionen}
Die Definition der trigonometrischen und hyperbolischen Funktionen 
sowie ihrer Umkehrfunktionen in einem Differentialkörper basieren
auf einer Exponential- und Logarithmus-Funktionen.
Da diese nicht eindeutig bestimmt sind, sind auch die
trigonometrischen und hyperbolischen Funktionen nicht eindeutig
bestimmt.
Die Exponentialfunktion ist nur bis auf einen konstanten Faktor bestimmt,
der sich natürlich auch in der Definition der trigonometrischen und
hyperbolischen Funktionen niederschlägt.
Dies scheint den Beziehungen
\[
\cos^2 x + \sin^2 x = 1
\qquad\text{und}\qquad
\cosh^2 x - \sinh^2 x = 1
\]
zu widersprechen, die für jede Wahl der Exponentialfunktion gelten.
Wir versuchen daher für die konkrete Wahl der Exponentialfunktion $e^x$
der klassischen Analysis den Einfluss eines Faktors auf die
trigonometrischen und hyperbolischen Funktionen zu ermitteln.
Im Folgenden reservieren wir daher die Notation $e^f$ und $\log f$
für die klassische definierten Exponential- und Logarithmusfunktionen.
Ebenso unterscheiden wir die verallgemeinerten trigonometrischen
und hyperbolischen Funktion mithilfe eines grossen Anfangsbuchstabens
von den klassischen Funtkionen.

Wir beginnen mit den hyperbolischen Funktionen.
Eine beliebige Exponentialfunktion $e(x)$ von $x$ ist ein Vielfaches
$a e^x$ der klassischen Exponentialfunktion.
Die hyperbolischen Funktionen sind dann
\begin{align*}
\operatorname{Cosh} x
&=
\frac{ae^x + e^{-x}/a}{2}
&&\text{und}&
\operatorname{Sinh} x
&=
\frac{ae^x - e^{-x}/a}{2}.
\intertext{Mit den Exponentialgesetzen folgt}
\operatorname{Cosh} x
&=
\frac{e^{x+\log a} + e^{-x-\log a}}{2}
&& \text{und}&
\operatorname{Sinh} x
&=
\frac{e^{x+\log a} - e^{-x-\log a}}{2}
\\
&=
\cosh(x+\log a)
&&&
&=
\sinh(x+\log a).
\end{align*}
Die verschiedenen möglichen hyperbolischen Funktionen entstehen daher
durch Verschiebung der unabhängigen Variablen.
Da man die Funktionen $\operatorname{Cosh}$ und $\operatorname{Sinh}$
als durch eine autonome Differentialgleichung definiert ansehen kann,
ist nachvollziehbar, dass die Definition die verschiedenen möglichen
Lösungen nicht unterscheiden kann.

Bei den trigonometrischen Funktionen betrachten wir einen Faktor der
Form $a=e^{i\delta}$, woraus $1/a = e^{-i\delta}$ folgt.
Die trigonometrischen Funktionen sind dann
\begin{align*}
\operatorname{Cos} x
&=
\frac12(ae^{ix}+e^{-ix}/a)
=
\frac12(e^{i\delta}e^{ix}+e^{ix}e^{-i\delta})
=
\frac12(e^{i(x+\delta)}+e^{-i(x+\delta)})
=
\cos(x+\delta)
\\
\operatorname{Sin} x
&=
\frac12(ae^{ix}-e^{-ix}/a)
=
\frac12(e^{i\delta}e^{ix}-e^{ix}e^{-i\delta})
=
\frac12(e^{i(x+\delta)}-e^{-i(x+\delta)})
=
\sin(x+\delta).
\end{align*}
Multiplikation mit einem Skalenfaktor bedeutet also wieder nur eine
Verschiebung des Arguments.

Die oben durchgeführten Rechnungen funktionieren natürlich auch für
einen komplexen Faktor, den wir $a=e^{b}$ schreiben können.
\begin{align*}
\operatorname{Cosh} x
&=
\frac12 ( ae^{x} + e^{-x}/a )
=
\frac12 ( e^{x+b} + e^{-x-b} )
=
\cosh(x+b)
\\
\operatorname{Sinh} x
&=
\frac12 ( ae^{x} - e^{-x}/a )
=
\frac12 ( e^{x+b} - e^{-x-b} )
=
\sinh(x+b).
\end{align*}
Entsprechende Formeln für die trigonometrischen Funktionen ergeben sich
daraus durch Multiplikation des Argumentes mit $i$.
In allen Fällen führen die verschiedenen Expontialfunktion zu einer
möglicherweise um einen komplexen Summanden verschobenen Funktion.

Bei den Umkehrfunktionen ist die Situation noch einfacher.
Sie sind als ein Logarithmus definiert.
Ein solcher ist nur bis auf einen konstanten Summanden definiert.
Diese Beobachtung ist auch konsistent mit der Interpretation der
Umkehrbeziehung für die Exponential- und Logarithmusfunktionen.
Die additive Konstante, um die sich die trigonometrischen und
hyperbolischen Umkehrfunktionen unterscheiden, ist auch die konstante
Translation, um die sich die trigonometrischen und hyperbolischen
Funktionen unterscheiden.

%
% Elementare Funktionen
%
\subsection{Elementare Funktionen und das Integrationsproblem}
Der Prozess der Erweiterung des Differentialkörpers der 
rationalen Funktionen $\mathbb{Q}(x)$ durch algebraische Elemente,
Logarithen und Exponentialfunktionen reproduziert also die Vorgehensweise
eines Mathematikers, der ein Integralproblem lösen will.
Er beginnt mit der Konstruktion eines Funktionenkörpers, der den
Integranden enthält.
Dazu wird es notwendig sein, zusätzliche Konstanten wie Wurzeln,
$\pi$, $i$ und andere oder Parameter hinzuzufügen.
Ausserdem müssen möglicherweise wiederholt Exponentialfunktionen,
Logarithmen und algebraische Funktionen hinzugefügt werden.
Alle Funktionen, mit denen ein menschlicher Rechner arbeitet, sind
auf diese Art konstruierbar.
Dies rechtfertigt, diesen Funktionen einen eigenen Namen zu geben.

\begin{definition}[Elementare Funktionen]
Sei $D$ eine Derivation in $\mathbb{Q}(x)$ mit $Dx=1$.
Eine Funktion $f(x)$ heisst \emph{elementar}, wenn $f$ in einer
Erweiterung von $\mathbb{Q}(x)$ enthalten ist, die durch hinzufügen
algebraischer Elemente, Logarithmen oder Exponentialfunktionen
entsteht.
\index{elementare Funktion}%
\end{definition}

Mit dem Begriff der elementaren Funktion lässt sich jetzt das
Problem der Integration in geschlossener Form präzis formulieren.
\index{geschlossene Form}%
Ein menschlicher Rechner erwartet, dass er den Differentialkörper,
in dem der Integrand gefunden werden kann, durch weitere Erweiterungen
so vergrössern kann, dass sich auch die Stammfunktion darin finden
lässt.
Die Stammfunktion ist daher auch eine elementare Funktion.
Das Problem der Integration in geschlossener Form ist daher das
folgende.

\begin{problem}[Integration in geschlossener Form]
\index{Integration in geschlossener Form}%
Gegebene eine elementare Funktion $f$ in einem Differentialkörper
$F$, finde eine elementare Funktion $\vartheta$ mit $D\vartheta = f$.
\end{problem}




